\subsection{Lensing Geometry and the Potential}\label{subsec:wl_geometry}

This section develops the geometry of weak gravitational lensing and introduces the lensing potential.

\subsubsection{Light bending by gravity}

Weak gravitational lensing is the small-angle bending of light rays by intervening masses. In the Born approximation, the bending angle is:

\begin{equation}
\alpha(\theta) = -2 \int dz \frac{dA(z)}{dA(z_s)} \frac{1}{2} \nabla_\perp \Phi(z)
\end{equation}

where $\Phi$ is the gravitational potential and $\nabla_\perp$ is the transverse Laplacian on the sky.

\subsubsection{The lensing potential}

For mathematical convenience, we define the lensing potential $\psi$ that encodes the cumulative bending:

\begin{equation}
\psi(\theta) = 2 \int_0^{z_s} dz \frac{D_A(z) D_A(z_s - z)}{D_A(z_s)} \Phi(z)
\end{equation}

where $D_A$ is angular diameter distance and $z_s$ is source redshift.

The convergence and shear are derived from $\psi$ via derivatives (see \cref{subsec:wl_convergence_shear}).

\remark{
The prefactor $D_A(z) D_A(z_s - z) / D_A(z_s)$ is the lensing efficiency: structures at intermediate distances bend rays most effectively, while structures near observer or source have reduced impact.
}

See \cref{subsec:background_distance} for angular diameter distance definitions.
